\chapter{进程}

\section{进程引入}
\concept{前趋图}
\begin{points}
  \item 前趋图是有向无循环图，用于描述程序段、语句或进程之间的先后执行关系。
  \item 若 $P_i \rightarrow P_j$，表示 $P_i$ 必须在 $P_j$ 开始之前完成。
  \item 没有前趋的结点称为初始结点，没有后继的结点称为终止结点。
  \item 做题时根据箭头判断哪些操作必须先完成，哪些操作可以并发。
\end{points}
\plain{用有向无环图画出“谁必须在谁之前完成”，目的是识别哪些操作能并发。}
\exam{常考给前趋关系画图、数可并发的语句、写同步信号量约束。}


\concept{程序的顺序执行}
\begin{points}
  \item 顺序执行指一个程序中的多个操作严格按照规定顺序执行，例如输入 $I \rightarrow$ 计算 $C \rightarrow$ 输出 $P$。
  \item 特征一：顺序性，每一步必须在下一步开始前完成。
  \item 特征二：封闭性，程序运行结果不受外界因素影响。
  \item 特征三：可再现性，只要初始条件和执行环境相同，每次运行结果相同。
\end{points}
\plain{一个程序自己按固定步骤走，外界不插手，所以结果可重复。}
\exam{常考顺序性、封闭性、可再现性三个特征。}


\concept{程序的并发执行}
\begin{points}
  \item 并发执行指多个程序或程序段在时间上重叠推进，一个程序尚未结束，另一个程序已经开始。
  \item 特征一：间断性，程序表现为“执行 - 暂停 - 执行”。
  \item 特征二：失去封闭性，多个程序共享资源，资源状态可能被他人改变。
  \item 特征三：不可再现性，在初始条件相同的情况下，执行顺序不同也可能导致结果不同。
\end{points}
\plain{并发执行这节重点不在和“并行”抠定义，而在看出它为什么会带来间断性、失去封闭性和结果不可再现。}
\exam{常考顺序执行和并发执行的三组特征对比，尤其为什么并发环境下会出现与时间有关的错误。}

\section{进程定义与描述}
\concept{进程的定义}
\begin{points}
  \item 进程是程序的一次执行过程，是系统进行资源分配和调度的基本单位之一。
  \item 进程不仅包含程序代码，还包含数据、堆栈、打开文件、寄存器现场、程序计数器和 PCB 等运行信息。
  \item 进程具有动态性、并发性、独立性、异步性和结构性。
\end{points}
\plain{这一定义里最关键的是“执行过程”和“资源分配调度单位”两层，不要把进程只理解成正在跑的代码片段。}
\exam{常考程序与进程区别、进程实体由程序段/数据段/PCB 构成、进程是资源分配基本单位。}


\concept{进程的特征}
\begin{points}
  \item \term{动态性}：进程是程序的一次执行过程，有创建、运行、暂停、撤销等生命周期；程序本身是静态实体。
  \item \term{并发性}：多个进程可同时存在于内存中，在一段时间内并发推进。
  \item \term{独立性}：进程是能独立运行、独立获得资源、独立接受调度的基本单位。
  \item \term{异步性}：进程按各自独立、不可预知的速度推进，因此执行结果可能受调度次序影响。
  \item \term{结构性}：进程实体通常由程序段、数据段和 PCB 等部分组成，也就是进程映像。
\end{points}
\plain{这五个特征里最容易混的是“动态性”和“并发性”：前者强调有生命周期，后者强调多个进程能重叠推进。}
\exam{简答题经常直接问“进程有哪些特征”，最好按这五点完整列出，而不是只答动态性和并发性。}


\concept{程序与进程的区别}
\begin{points}
  \item 程序是静态的代码文件，进程是动态的执行活动。
  \item 一个程序可以对应多个进程，例如多个用户同时打开同一个编辑器。
  \item 一个进程在运行中可能执行多个程序片段，例如通过 \code{exec} 装入新程序。
  \item 程序没有生命周期状态，进程有创建、就绪、运行、阻塞、终止等状态。
\end{points}
\plain{区分程序和进程时，最稳的是按“静态/动态、是否拥有生命周期、是否携带资源和状态”三条线来答。}
\exam{常考判断题：同一程序能否对应多个进程、进程能否在运行中装入新程序、程序本身是否具有状态。}


\concept{PCB 进程控制块}
\begin{points}
  \item PCB 是操作系统感知和管理进程的唯一标志。
  \item PCB 保存进程标识信息、处理机状态、调度信息、进程控制信息、内存管理信息、I/O 状态信息等。
  \item 上下文切换本质上是保存当前进程 PCB 中的现场，并从下一个进程 PCB 中恢复现场。
  \item 选择题直觉：问“进程存在的唯一标志”通常答 PCB，不是程序、PID 或页表。
\end{points}
\plain{PCB 就是 OS 给每个进程建的档案，没有 PCB，OS 就没法调度和恢复这个进程。}
\exam{常考 PCB 内容、PCB 的作用、进程存在的唯一标志。}


\concept{PCB 的组织方式}
\begin{points}
  \item 为了便于创建、撤销、阻塞、唤醒和调度进程，操作系统需要把大量 PCB 组织起来。
  \item \term{线性方式}：把 PCB 顺序存放在连续内存中，实现简单，但插入删除不够灵活。
  \item \term{链接方式}：把同一状态的 PCB 按链表连接，便于按状态组织就绪队列、阻塞队列。
  \item \term{索引方式}：建立索引表，再由索引项指向对应 PCB，便于按标识快速定位。
\end{points}
\plain{PCB 不只是一个个单独存在，OS 还得把它们按状态串起来或索引起来，否则调度器没法高效找到“下一个该处理谁”。}
\exam{PPT 专门讲了 PCB 组织方式，常考线性/链接/索引三种方法的适用特点，尤其链表为什么适合状态队列。}


\concept{Bernstein 条件}
\begin{points}
  \item 为判断两个程序段能否并发执行且不产生与时间有关的错误，可引入读集 $R(S_i)$ 和写集 $W(S_i)$。
  \item 若满足 $R(S_i)\cap W(S_j)=\varnothing$，$R(S_j)\cap W(S_i)=\varnothing$，$W(S_i)\cap W(S_j)=\varnothing$，则 $S_i$ 与 $S_j$ 可安全并发。
  \item 前两个条件避免“读到被对方修改中的数据”，第三个条件避免“两个程序段写同一数据导致结果丢失”。
  \item 满足这三个条件时，并发执行仍可保持封闭性和可再现性。
\end{points}
\plain{Bernstein 条件本质是在问：两段程序会不会互相读坏、写坏对方要用的数据。}
\exam{常考给出若干语句的读集写集，判断哪几条可以并发执行。}


\concept{\code{cobegin} 与 \code{coend} 描述并发}
\begin{points}
  \item 并发语句常写作 \code{cobegin S1; S2; ...; Sn; coend}。
  \item 其含义是多个语句段在逻辑上可并发执行，但 \code{coend} 之后的语句必须等这些并发段都结束后才能继续。
  \item 它常与前趋图一起使用，用于描述程序并发结构。
\end{points}
\plain{\code{cobegin} 像“这些分支一起干”，\code{coend} 像“所有分支会合后再往下走”。}
\exam{常考把前趋图改写成 \code{cobegin/coend} 形式，或反过来画前趋关系。}

\concept{为什么引入进程}
\begin{points}
  \item 引入进程是为了刻画程序执行过程的动态性，而不再只把程序看成静态指令集合。
  \item 引入进程有助于发挥系统并发性，提高资源利用率。
  \item 引入进程还用于处理共享性问题，正确描述程序执行状态及系统对运行实体的管理。
\end{points}
\plain{只有“程序”这个静态概念时，很难描述并发推进、资源共享和状态变化，所以 OS 需要再引入“进程”这个动态执行实体。}
\exam{若问“操作系统为什么引入进程”，答题可按“刻画动态性、支持并发、处理共享与状态管理”三点展开。}


\concept{进程映像}
\begin{points}
  \item 进程映像是某时刻进程的内容及其状态的集合，也可理解为进程在系统中的完整存在形态。
  \item 通常包括 PCB，以及进程在虚拟地址空间中的代码段、数据段、栈、堆等内容。
  \item 某些教材还会强调核心栈，用于保存进程在核心态工作时的现场保护。
\end{points}
\plain{进程映像回答的是“这个进程完整长什么样”，不只是代码，还包括数据和管理信息。}
\exam{常考进程映像由哪些部分组成，以及它和程序本身的区别。}


\concept{进程上下文}
\begin{points}
  \item 进程物理实体及其支持运行的环境合称为进程上下文，进程运行可看作在其上下文中执行。
  \item 进程映像强调“全部内容集合”；上下文更强调“当前运行现场和支撑环境”。
  \item 当系统调度新进程占有处理器时，就会发生上下文切换。
  \item 上下文可分为用户级上下文、寄存器级上下文和系统级上下文。
\end{points}
\plain{进程映像偏“全家福”，进程上下文偏“此刻跑到哪、现场是什么样”。}
\exam{常考进程映像和进程上下文的区别，以及上下文切换时保存恢复什么。}


\concept{进程上下文的组成}
\begin{points}
  \item \term{用户级上下文}：程序正文、数据、共享存储区和用户栈等，属于进程虚拟地址空间中的用户态部分。
  \item \term{寄存器级上下文}：程序计数器、程序状态字、栈指针、通用寄存器、控制寄存器等现场信息。
  \item \term{系统级上下文}：PCB、内核栈、页表等主存管理信息、I/O 状态、文件描述符表等内核管理信息。
  \item 上下文切换时，并不是只换一组寄存器，而是按需要切换这些和运行现场直接相关的内容。
\end{points}
\plain{用户级上下文偏“进程自己的地址空间内容”，寄存器级上下文偏“CPU 此刻现场”，系统级上下文偏“内核为这个进程维护的管理环境”。}
\exam{这类题常考“上下文切换保存什么”，按用户级/寄存器级/系统级三层去答最稳。}


\section{进程状态与转换}
\concept{三状态模型}
\begin{points}
  \item \term{就绪态}：进程已具备运行条件，只差 CPU。
  \item \term{运行态}：进程正在 CPU 上执行。
  \item \term{阻塞态/等待态}：进程因等待某事件不能运行，即使 CPU 空闲也不能执行。
  \item 状态转换：就绪 $\rightarrow$ 运行由调度引起；运行 $\rightarrow$ 就绪由时间片到或被抢占引起；运行 $\rightarrow$ 阻塞由等待 I/O/资源引起；阻塞 $\rightarrow$ 就绪由等待事件完成引起。
\end{points}
\plain{三态模型先抓最核心的三种活动状态：能不能运行、正在不在运行、是不是在等事件。}
\exam{常考一个处理器上运行态进程最多几个，以及三态之间的转换条件。}


\concept{五状态模型}
\begin{points}
  \item 在三状态基础上加入新建态和终止态。
  \item 新建态表示进程正在创建，尚未进入就绪队列。
  \item 终止态表示进程执行结束或被撤销，等待系统回收资源。
  \item 创建进程通常涉及分配 PCB、建立地址空间、装入程序、初始化资源、进入就绪队列。
\end{points}
\plain{五态就是在三态前后各补一个“进来”和“出去”的阶段，把完整生命周期补齐。}
\exam{常考五态模型比三态多了哪两个状态，以及创建、撤销分别对应什么阶段。}


\concept{挂起状态}
\begin{points}
  \item 当内存紧张或系统需要调节多道程序度时，可把进程换出到外存，形成挂起状态。
  \item 就绪挂起：进程不在内存，但一旦调入内存即可运行。
  \item 阻塞挂起：进程不在内存，且还在等待某事件。
  \item 中级调度/交换调度常与挂起状态相关。
\end{points}
\plain{挂起强调“暂时不在内存里”，它和阻塞不是一回事；一个进程既可能阻塞，也可能同时被挂起。}
\exam{常考为什么引入挂起状态，以及就绪挂起、阻塞挂起的区别。}

\concept{挂起进程的特征}
\begin{points}
  \item 挂起进程可视为暂时不在主存中的进程，在被重新调入主存前不参与处理机调度。
  \item 挂起进程不能立即执行，即使它原来已经具备运行条件也不行。
  \item 挂起进程可能仍在等待某个事件，但该事件结束并不自动解除挂起状态。
  \item 进入挂起常由操作系统、父进程或进程自身发起；结束挂起通常需由操作系统或父进程发出激活命令。
\end{points}
\plain{判断“挂起”时，关键看它是否已经被换出主存、被暂时禁止参与调度，而不只是看它有没有在等事件。}
\exam{区分“等待”和“挂起”时，要答出：等待针对事件，挂起针对是否驻留主存以及是否允许参与调度。}


\section{进程组织、控制与管理}
\concept{进程队列}
\begin{points}
  \item 就绪队列：所有驻留内存且处于就绪状态、等待 CPU 的进程集合。
  \item 设备队列：等待某个 I/O 设备的进程集合。
  \item 作业队列：外存上等待进入内存的作业集合。
  \item 进程在生命周期中会在多个队列之间迁移，调度程序就是从这些队列中选择合适对象。
\end{points}
\plain{队列是 OS 组织进程的方式：就绪的排就绪队列，等设备的排设备队列。}
\exam{常考进程在生命周期中如何在不同队列迁移。}


\concept{进程控制原语}
\begin{points}
  \item 进程控制的常见功能包括创建、撤销、阻塞、唤醒、挂起和激活。
  \item 这些功能通常由内核原语实现，原语由若干机器指令构成，执行期间不可分割。
  \item 原语必须保证原子性，否则在修改 PCB、状态和队列时可能被中断，导致系统状态不一致。
\end{points}
\plain{进程控制题先抓住“原语”二字：它是内核里不可分割的基本操作，专门负责改 PCB、改状态、改队列。}
\exam{常考什么叫原语、为什么原语必须原子执行，以及哪些进程管理操作属于控制原语。}


\concept{创建原语}
\begin{points}
  \item 创建进程需要申请唯一标识、分配 PCB、分配必要资源、初始化上下文、插入就绪队列。
  \item 引起创建的常见原因包括系统初始化、用户请求和父进程通过系统调用创建子进程。
  \item 创建后的父子进程可以并发执行，也可以由父进程等待子进程结束后再继续。
  \item 为描述父子关系，可用进程树/进程图表示“谁创建了谁”。
  \item Unix/Linux 中常见系统调用：\code{fork()} 创建子进程，\code{exec()} 装入新程序，\code{wait()} 等待子进程，\code{exit()} 退出进程。
  \item \code{fork()} 返回两次：父进程中返回子进程 PID，子进程中返回 0；失败返回负值。
\end{points}
\plain{创建原语的标准流程就是“要 PCB、配资源、填初值、进就绪队列”，代码题再把它和 \code{fork/exec/wait} 对起来。}
\exam{常考创建原语步骤、进程树、fork/exec/wait/exit，以及多次 \code{fork()} 后的进程数和输出行数。}


\concept{撤销原语}
\begin{points}
  \item 进程撤销的原因可分为正常结束、异常结束和外界干预。
  \item 异常结束常见于越界、非法指令、I/O 故障、内存不足等；外界干预常见于父进程或系统强制终止。
  \item 撤销原语的主要工作是找到对应 PCB，必要时先停止其执行，再回收该进程占有的资源和 PCB。
  \item 题目还可能区分两种策略：只撤销指定进程，或连同其全部子孙进程一起撤销。
\end{points}
\plain{撤销原语不是简单“结束程序”，而是把进程从系统里彻底清掉，包括停执行、收资源、删 PCB。}
\exam{常考进程终止原因分类、撤销原语步骤，以及“撤销某进程是否连带其子孙进程”这类策略题。}


\concept{阻塞原语与唤醒原语}
\begin{points}
  \item 进程因等待系统服务、I/O 完成、合作进程配合或某种事件到来时，可由执行态转为阻塞态。
  \item 阻塞原语的核心动作是停止当前进程执行、保存 CPU 现场、把状态改为阻塞，并插入相应等待队列。
  \item 当等待事件发生后，唤醒原语将该进程移出等待队列，改为就绪态，再插入就绪队列。
  \item 阻塞通常由进程自己触发；唤醒通常由发现事件发生的其他进程或内核触发。
\end{points}
\plain{阻塞和唤醒一定成对理解：自己因为等事件而睡下去，别人或内核在事件到了之后把你放回就绪队列。}
\exam{常考阻塞/唤醒的触发条件、两类原语各自修改哪种队列，以及“阻塞态不能直接变运行态”。}


\concept{挂起原语与激活原语}
\begin{points}
  \item 挂起强调把进程暂时从内存换出到外存，常用于缓解内存压力或调整多道程序度。
  \item 挂起时，活动就绪可变为静止就绪，活动阻塞可变为静止阻塞；若进程正在执行，还需先保护现场并重新调度。
  \item 激活原语的作用与之相反：把挂起进程重新调入内存，使其回到就绪或阻塞的活动状态。
  \item 挂起/激活既可针对指定进程，也可能连同其子孙进程一起处理。
\end{points}
\plain{挂起和阻塞别混：阻塞是“等事件”，挂起是“先不放在内存里”；一个进程可以既阻塞又挂起。}
\exam{常考活动就绪/静止就绪、活动阻塞/静止阻塞之间如何转换，以及挂起与激活对状态和调度的影响。}


\concept{进程切换与模式切换}
\begin{points}
  \item 模式切换是用户态和内核态之间的切换，不一定更换正在运行的进程。
  \item 进程切换是 CPU 从一个进程切换到另一个进程，一定需要内核介入，通常伴随上下文保存和恢复。
  \item 系统调用会导致模式切换；是否导致进程切换，要看是否阻塞或调度。
  \item I/O 请求通常先从用户态陷入内核态，若需要等待 I/O 完成，则当前进程阻塞并发生进程切换。
\end{points}
\plain{模式切换只是换权限；进程切换还要保存/恢复另一个进程的现场，开销更大。}
\exam{常考系统调用是否必然进程切换、进程切换需要保存哪些上下文。}


\section{进程通信}
\concept{进程协作的原因}
\begin{points}
  \item 信息共享：多个进程需要交换或共享感兴趣的数据。
  \item 提高运算速度：把任务分解后并发执行，以缩短总完成时间。
  \item 系统模块化：把复杂系统拆成多个相互协作的模块或进程。
\end{points}
\plain{进程通信存在的意义，不只是“把消息传过去”，而是为了共享信息、并发加速以及支撑系统模块化设计。}
\exam{若问“为什么需要进程通信”，不要只答数据交换，还要补上并发加速和系统模块化。}

\concept{低级通信与高级通信}
\begin{points}
  \item 低级通信主要通过信号量等同步工具传递少量控制信息。
  \item 高级通信用于传递大量数据或结构化消息，常见方式包括共享存储、消息传递、管道。
\end{points}
\plain{共享存储的优势是“数据直接放一块公共内存里”，所以传大块数据最快；代价是同步互斥必须自己另外补上。}
\exam{常考为什么共享存储速度快、为什么它又最容易和信号量/互斥锁联动出题。}

\concept{进程通信的基本形态}
\begin{points}
  \item 最基本的两类 IPC 模型是消息传递模型和共享内存模型。
  \item 消息传递模型通过发送和接收消息交换信息，进程之间不必直接共享地址空间。
  \item 共享内存模型让多个进程直接访问同一片内存区域，数据交换效率高，但必须额外处理同步互斥。
\end{points}
\plain{IPC 题先抓总类：要么靠 OS 传消息，要么大家直接看同一块共享内存。邮箱、消息队列、管道、共享存储，基本都能往这两类里归。}
\exam{若题目先问“IPC 的两种基本模型是什么”，就先答消息传递和共享内存，再展开各自优缺点。}


\concept{共享存储}
\begin{points}
  \item 多个进程共享同一块内存区域，通过读写共享区交换数据。
  \item 优点：速度快，适合大量数据交换。
  \item 缺点：必须另行解决同步互斥问题，防止同时读写导致不一致。
\end{points}
\plain{共享存储传数据最快，但同步问题也最直接暴露出来，所以它常和信号量、互斥一起考。}
\exam{常考共享存储为什么速度快、为什么还必须配合同步机制。}


\concept{消息传递}
\begin{points}
  \item 进程通过发送和接收消息通信，可由操作系统提供消息队列、邮箱等机制。
  \item 直接通信：发送者明确指定接收者。
  \item 间接通信：通过邮箱或消息队列进行。
  \item 优点：同步控制相对清晰，适合分布式系统；缺点是消息复制和内核介入可能带来开销。
\end{points}
\plain{消息传递的核心是由操作系统充当中介，发送方和接收方不必直接共享地址空间，但要付出内核介入和消息复制的代价。}
\exam{常考直接通信 vs 间接通信、阻塞 vs 非阻塞接收发送，以及消息传递为什么更适合分布式环境。}


\concept{管道}
\begin{points}
  \item 管道是一种半双工通信机制，常用于具有亲缘关系的进程之间。
  \item 匿名管道常用于父子进程；命名管道可用于无亲缘关系的进程。
  \item Shell 中的 \code{|} 就是把前一个命令的输出作为后一个命令的输入。
\end{points}
\plain{管道更像一根有容量限制的字节流通道，题目里除了“半双工”外，还常顺带考互斥、同步和写满/读空时谁该阻塞。}
\exam{常考匿名管道和命名管道区别、为什么匿名管道常用于亲缘进程、以及双向通信为何通常要两根管道。}
